\section{Introduction}
\label{sec:introduction}

Sixth-generation (6G) wireless networks are expected to provide ubiquitous
connectivity for massive Internet of Things (IoT) deployments
\cite{cai2026,xu2026}. To serve devices in remote areas with limited
terrestrial coverage, low Earth orbit (LEO) satellites complement cellular
networks for applications such as smart agriculture and climate monitoring
\cite{shen2025,kodheli2021}. These power-limited devices transmit short
packets sporadically, with only a small fraction active at any given time
\cite{cai2026,shen2025,mei2022}. In grant-based random access,
pre-transmission handshaking incurs substantial signaling overhead, while the
long propagation delay of terrestrial--satellite links further increases
access latency \cite{cai2026,xu2026}. To reduce latency and signaling
overhead, grant-free random access (GFRA) allows active devices to transmit
pilots and data directly without waiting for a scheduling grant
\cite{cai2026,xu2026}.

Reliable communication over LEO links requires handling large delays and
Doppler shifts caused by long propagation distances and high satellite
mobility. The 3rd Generation Partnership Project (3GPP) Release 17 framework
addresses these effects through device-side precompensation assisted by the
global navigation satellite system (GNSS) \cite{3gpp38300}. With GNSS
assistance, active devices can align their uplink transmissions so that
signals arriving at the satellite receiver exhibit only small residual delays
and Doppler shifts \cite{shen2025,li2023}. However, low-cost IoT devices may
not be equipped with GNSS receivers or may fail to obtain valid positioning
due to signal blockage or interference. To support satellite access under
these conditions, the 3GPP Release 19 study specifically considers
GNSS-independent operation \cite{3gpp22865}. Without device-side
precompensation, the satellite receiver must accommodate large differential
delays and Doppler shifts, which cause severe inter-symbol and inter-carrier
interference in conventional orthogonal frequency division multiplexing (OFDM)
systems \cite{shen2025,huang2026}.

To mitigate such doubly dispersive effects, delay--Doppler and chirp-based
modulation schemes have been investigated. Orthogonal time frequency space
(OTFS) modulation exploits the quasi-static channel representation in the
delay--Doppler domain \cite{shen2025,hadani2017}. Alternatively, affine
frequency division multiplexing (AFDM) provides a low-complexity modulation
structure based on one-dimensional chirp signals \cite{huang2026,ding2026}.
Its modulator can be implemented using a single inverse fast Fourier transform
(IFFT) and two stages of pointwise chirp multiplication \cite{bemani2022}.
Moreover, its chirp parameters can be configured according to the delay and
Doppler spreads, making AFDM a promising candidate for satellite access
\cite{bemani2022}.

Since device activity states and channels are unknown at the receiver, GFRA
requires joint activity detection and channel estimation (JADCE)
\cite{cai2026,shen2025,zhang2020_brmp}. Existing methods exploit sporadic
device activity to formulate JADCE as a compressed sensing (CS) problem
\cite{liu2018,zhang2019_dnn_bsbl}, incorporating structured priors to capture
channel sparsity \cite{shen2025,ding2026,shen2022_mimootfs,ke2020} and spatial
correlation \cite{djelouat2024_hierarchical}. For AFDM systems, a graph neural
network (GNN) was developed in \cite{ding2026} to adaptively learn sparsity
patterns, though the channel remains represented on a predefined
delay--Doppler grid.

In practical propagation environments, however, power leakage caused by
fractional delays and Doppler shifts can weaken the sparsity of the AFDM
effective channel matrix \cite{huang2026}. While off-grid sparse Bayesian
learning (SBL) and grid-refinement methods account for fractional delay or
Doppler parameters in OTFS \cite{WeiEtAl2022OffgridSBL} and AFDM
\cite{li2025_afdm_sbl,luo2026_afdm_fdfd}, they still rely on a sparse path
representation. This motivates data-driven channel estimation using learned
denoisers \cite{metzler2016,he2018_ldamp} or generative priors
\cite{balevi2021_generative} to capture channel structures beyond the assumed
sparse model \cite{cai2026,huang2026}.

Score-based generative models learn channel priors directly from samples for
channel estimation
\cite{song2019,arvinte2023,zhou2025_diffusion,chen2025_variational,fesl2024_diffusion}.
For AFDM, score-based learning has been applied to the effective channel to
address off-grid leakage \cite{huang2026}. For massive random access, a
two-stage receiver was proposed in \cite{ying2026} that first identifies
active devices via a Transformer-based detector and then performs joint
channel estimation and data detection using diffusion models. However, both
approaches rely on iterative posterior sampling, which requires repeated
score-function evaluations and incurs high computational complexity
\cite{huang2026,ying2026}.

To avoid iterative posterior sampling, the score-based turbo message passing
(STMP) algorithm incorporates generative channel priors into JADCE for OFDM
systems \cite{cai2026}. However, when applied to finite-dimensional AFDM
systems, its Gaussian approximation of the channel posterior may not accurately
capture the activity likelihoods associated with structured effective channels.
Its use of separate score networks for channel mean and variance estimation
also increases model storage requirements for onboard deployment.

Against this background, we investigate GNSS-independent AFDM GFRA for LEO
satellite IoT systems. Our main contributions are summarized as follows:

\begin{enumerate}[label=\arabic*)]

\item We propose a multi-frame AFDM transmission scheme to handle large differential
delays and Doppler shifts at the satellite receiver without device-side
precompensation. By deriving the input--output relationship, we show that the
effective channels across frames differ only by Doppler-induced phase
rotations. Consequently, transmitting additional pilot frames provides more
observations without increasing the number of unknown channel coefficients.

\item We develop a message-passing receiver for JADCE that iterates between
multi-frame linear estimation (MFLE) for inter-user interference suppression
and generative-prior-aided joint denoising (GPJD). To handle fractional delays
and Doppler shifts without sparse or grid-aligned assumptions, GPJD employs a
score-based generative prior learned from AFDM effective-channel samples.
Exploiting the analytical relationship between the posterior mean and
variance, the same denoiser provides both channel and variance estimates
without a separately trained variance model. Moreover, activity probabilities
are updated using a density-ratio estimator trained via binary classification,
avoiding a Gaussian approximation when evaluating the activity likelihood ratio.

\item Numerical results show that the proposed receiver outperforms the adapted STMP
baseline in both activity detection and channel estimation, requiring fewer
pilot frames to achieve the same estimation accuracy. The results also confirm
the effectiveness of the density-ratio activity update, the robustness of
pilot-based phase estimation, and strong generalization across non-terrestrial
network tapped delay line (NTN-TDL) channel profiles without retraining.

\end{enumerate}

The remainder of this paper is organized as follows. Section~\ref{sec:system}
presents the system model and the multi-frame AFDM input--output relationship.
Sections~\ref{sec:algorithm}--\ref{sec:overall} describe the message-passing
framework, the learned channel denoiser and activity classifier, and the
overall algorithm with its computational complexity analysis, respectively.
Section~\ref{sec:simulations} presents the simulation results, and
Section~\ref{sec:conclusion} concludes this paper.
